% assets/clase13-runge-lineal.tex
% Función de Runge f(x)=1/(1+25x^2) en [-1,1] vs. su interpolante lineal a
% trozos con 9 nodos equiespaciados. Ilustra §3.4 del summary: a diferencia
% de la interpolación global de grado alto (Clase 12), la lineal a trozos
% no oscila cerca de los bordes.
\begin{figure}[H]
  \centering
  \begin{tikzpicture}
    \begin{axis}[interpfig,
      xlabel={$x$}, ylabel={$y$},
      xmin=-1.05, xmax=1.05, ymin=-0.05, ymax=1.12,
      % el eje x va abajo (no centrado): rótulos en la punta de cada eje
      xlabel style={at={(ticklabel* cs:1)}, anchor=west},
      ylabel style={at={(ticklabel* cs:1)}, anchor=south},
      % leyenda debajo del eje, fuera de la curva
      legend style={at={(0.5,-0.2)}, anchor=north}, legend columns=2,
      height=5.2cm,
    ]
      \addplot[figblue, line width=1.1pt, domain=-1:1, samples=200, forget plot]
        {1/(1+25*x^2)};
      \addlegendimage{figblue, line width=1.1pt}
      \addlegendentry{$f(x)=1/(1+25x^2)$}

      \addplot[figred, line width=1pt, mark=none]
        coordinates {
          (-1,0.038462) (-0.75,0.066390) (-0.5,0.137931) (-0.25,0.390244)
          (0,1) (0.25,0.390244) (0.5,0.137931) (0.75,0.066390) (1,0.038462)
        };
      \addlegendentry{lineal a trozos ($n=8$)}

      \addplot[only marks, mark=*, mark size=2pt, mark options={fill=figamber, draw=figblue}, forget plot]
        coordinates {
          (-1,0.038462) (-0.75,0.066390) (-0.5,0.137931) (-0.25,0.390244)
          (0,1) (0.25,0.390244) (0.5,0.137931) (0.75,0.066390) (1,0.038462)
        };
    \end{axis}
  \end{tikzpicture}
  \caption{Función de Runge y su interpolante lineal a trozos con 9 nodos equiespaciados. A diferencia del interpolante polinomial global (Clase 12), no hay oscilación cerca de los bordes.}
  \label{fig:clase13-runge-lineal}
\end{figure}
